\chapter{Linear Algebra for State Estimation}
\label{ch:linalg}

\begin{objectives}
\begin{itemize}
  \item Master the cross-product matrix $\sx{a}$ and the identities used in every Jacobian derivation.
  \item Understand the four fundamental subspaces, especially the \emph{left nullspace} on which the
        MSCKF rests.
  \item Know when and why to use Cholesky, LDL$^\mathsf{T}$, QR and SVD, and why a filter should never
        form an explicit inverse.
  \item Use block matrices, Schur complements and the Woodbury identity fluently.
  \item Understand the matrix exponential and floating-point arithmetic well enough to write
        numerically robust code.
\end{itemize}
\end{objectives}

Most of the computation in a Kalman filter is linear algebra: products of covariance matrices,
solutions of linear systems, and factorizations. This chapter collects the tools the rest of the
notes rely on and explains the numerical reasoning behind them. Readers fluent in numerical linear
algebra can skim it, but should not skip \cref{sec:skew,sec:subspaces,sec:floating}.

% ------------------------------------------------------------------------------
\section{Vectors, coordinates and the cross product}
\label{sec:skew}

\subsection{Vectors versus their coordinates}

A physical vector, such as a velocity or the direction of gravity, exists independently of any
coordinate system. To compute with it we express it in a \emph{frame}: an origin plus three
orthonormal, right-handed axes. The same physical vector $\vec{v}$ has different coordinate
triples $v_w \in \R^3$ and $v_b \in \R^3$ in the world and body frames, and the two are related by
a rotation matrix, $v_w = R_{wb}\,v_b$ (\cref{ch:rotations}). Confusing a vector with its
coordinates in some frame is the root of most frame bugs. Throughout these notes, a subscript on a
coordinate vector names the frame it is expressed in.

For coordinate vectors $a, b \in \R^3$ we use the standard inner product
$a\T b = \sum_i a_i b_i$, the Euclidean norm $\|a\| = \sqrt{a\T a}$, and the cross product
$a \times b$. The cross product is bilinear, anti-commutative ($a \times b = -b \times a$), and
orthogonal to both arguments. Its norm is $\|a\|\|b\|\sin\angle(a,b)$.

\subsection{The cross-product (skew-symmetric) matrix}

Because $a \times b$ is linear in $b$, it can be written as a matrix-vector product.

\begin{definition}[Skew-symmetric matrix]
For $a = (a_1, a_2, a_3)\T$, the \emph{cross-product matrix} of $a$ is
\begin{equation}
  \sx{a} = \begin{bmatrix} 0 & -a_3 & a_2 \\ a_3 & 0 & -a_1 \\ -a_2 & a_1 & 0 \end{bmatrix},
  \qquad\text{so that}\qquad \sx{a}\, b = a \times b \quad \forall\, b \in \R^3 .
  \label{eq:skew}
\end{equation}
It is skew-symmetric: $\sx{a}\T = -\sx{a}$. Conversely, every $3\times3$ skew-symmetric matrix
$S$ equals $\sx{s}$ for the unique vector $s = (S_{32}, S_{13}, S_{21})\T$. We write
$s = S\veeop$ for this inverse operation, the \emph{vee} map. In Lie theory, $\sx{\cdot}$ is
also written $(\cdot)\hatop$ and called the \emph{hat} map.
\end{definition}

The following identities are used constantly in derivations. Learn to recognize them on sight.

\begin{proposition}[Cross-product identities]
\label{prop:skew-identities}
For $a, b, c \in \R^3$ and $R \in \SO$ (a rotation matrix):
\begin{align}
  \sx{a}\, b &= -\sx{b}\, a, \label{eq:skew-anticomm}\\
  \sx{a}\, a &= 0, \label{eq:skew-self}\\
  \sx{a}^2 &= a a\T - \|a\|^2 I, \label{eq:skew-sq}\\
  \sx{a}^3 &= -\|a\|^2 \sx{a}, \label{eq:skew-cube}\\
  \sx{a}\sx{b} &= b a\T - (a\T b)\, I, \label{eq:skew-prod}\\
  \sx{a \times b} &= \sx{a}\sx{b} - \sx{b}\sx{a} = b a\T - a b\T, \label{eq:skew-commutator}\\
  \sx{R a} &= R \sx{a} R\T, \label{eq:skew-rot}\\
  R\,(a \times b) &= (Ra) \times (Rb). \label{eq:cross-rot}
\end{align}
\end{proposition}

\begin{proof}
\eqref{eq:skew-anticomm} and \eqref{eq:skew-self} restate $a\times b = -b\times a$ and
$a \times a = 0$. For \eqref{eq:skew-prod}, apply both sides to an arbitrary $c$ and use the
``BAC--CAB'' rule $a \times (b \times c) = b\,(a\T c) - c\,(a\T b)$:
$\sx{a}\sx{b}c = a\times(b\times c) = b a\T c - (a\T b) c$. Setting $b = a$ gives
\eqref{eq:skew-sq}, and multiplying \eqref{eq:skew-sq} by $\sx{a}$ and using \eqref{eq:skew-self}
gives \eqref{eq:skew-cube}. Subtracting \eqref{eq:skew-prod} from its version with $a, b$ swapped
gives the second equality of \eqref{eq:skew-commutator}. The first follows from the Jacobi identity
$a \times (b \times c) + b \times (c \times a) + c \times (a \times b) = 0$. For
\eqref{eq:cross-rot}, the cross product is defined by orientation and lengths, both preserved by a
proper rotation. Finally \eqref{eq:skew-rot} follows from \eqref{eq:cross-rot}:
$\sx{Ra}\, c = (Ra)\times c = (Ra)\times(R R\T c) = R\big(a \times R\T c\big) = R\sx{a}R\T c$.
\end{proof}

\begin{keyidea}[moving a vector out of a cross product]
Identity \eqref{eq:skew-anticomm} is how unknowns are isolated in Jacobian derivations. A
perturbation appears as $\sx{\delta\theta}\,v$, and the Jacobian with respect to $\delta\theta$ is
read off after rewriting it as $-\sx{v}\,\delta\theta$.
\end{keyidea}

Identity \eqref{eq:skew-cube} is what makes closed-form rotation formulas possible: every power of
$\sx{a}$ reduces to a multiple of $\sx{a}$ or $\sx{a}^2$, so infinite series in $\sx{a}$ collapse
to two terms (\cref{ch:rotations}).

% ------------------------------------------------------------------------------
\section{The four fundamental subspaces}
\label{sec:subspaces}

Let $A \in \R^{m\times n}$ have rank $r$. Four subspaces are associated with it:
\begin{center}
\begin{tabular}{@{}l l l l@{}}
\toprule
Subspace & Definition & Lives in & Dimension \\
\midrule
column space (range) & $\range(A) = \{Ax\}$ & $\R^m$ & $r$ \\
nullspace & $\nullsp(A) = \{x : Ax = 0\}$ & $\R^n$ & $n - r$ \\
row space & $\range(A\T)$ & $\R^n$ & $r$ \\
left nullspace & $\nullsp(A\T) = \{y : y\T A = 0\}$ & $\R^m$ & $m - r$ \\
\bottomrule
\end{tabular}
\end{center}

\begin{theorem}[Fundamental theorem of linear algebra]
\label{thm:fundamental}
$\nullsp(A) = \range(A\T)^\perp$ and $\nullsp(A\T) = \range(A)^\perp$. Consequently
$\R^n = \range(A\T) \oplus \nullsp(A)$ and $\R^m = \range(A) \oplus \nullsp(A\T)$, with
orthogonal direct sums.
\end{theorem}
\begin{proof}
$y \in \nullsp(A\T) \iff A\T y = 0 \iff x\T A\T y = 0\ \forall x \iff (Ax)\T y = 0\ \forall x
\iff y \perp \range(A)$. The first statement is the same argument applied to $A\T$.
\end{proof}

The \emph{left nullspace} deserves special attention because the MSCKF is built on it. Suppose a
stacked measurement depends on an unwanted variable $f$ through a tall matrix $H_f \in \R^{m\times
3}$ of full column rank:
\begin{equation}
  r = H_x\,\delta x + H_f\,\delta f + n .
\end{equation}
Take any matrix $A \in \R^{m\times(m-3)}$ whose columns form an orthonormal basis of
$\nullsp(H_f\T)$. Then $A\T H_f = 0$, and multiplying through by $A\T$ gives
$A\T r = A\T H_x\,\delta x + A\T n$: the unwanted variable has disappeared. \Cref{ch:msckfupdate}
shows that no information about $\delta x$ is lost in the process.

\begin{definition}[Orthogonal projector]
A matrix $\Pi$ is an orthogonal projector if $\Pi^2 = \Pi$ and $\Pi\T = \Pi$. If the columns of
$Q$ are an orthonormal basis of a subspace $\mathcal{S}$, the projector onto $\mathcal{S}$ is
$\Pi_{\mathcal{S}} = Q Q\T$. If $A$ has full column rank, the projector onto $\range(A)$ is
$A (A\T A)\inv A\T$, and $I - \Pi_{\mathcal{S}}$ projects onto $\mathcal{S}^\perp$.
\end{definition}

For the matrix $A$ spanning the left nullspace of $H_f$, these facts combine into an identity we
will use to prove the MSCKF's optimality:
\begin{equation}
  A A\T = I - H_f (H_f\T H_f)\inv H_f\T .
  \label{eq:nullspace-projector}
\end{equation}

% ------------------------------------------------------------------------------
\section{Symmetric and positive-definite matrices}

Covariance matrices are symmetric and positive semi-definite, and much of the numerical care in a
filter goes into keeping them that way.

\begin{definition}
A symmetric matrix $P = P\T \in \R^{n\times n}$ is \emph{positive semi-definite} (PSD, $P \succeq 0$) if
$x\T P x \ge 0$ for all $x$, and \emph{positive definite} (PD, $P \succ 0$) if $x\T P x > 0$
for all $x \ne 0$. We write $P_1 \succeq P_2$ (the Loewner order) if $P_1 - P_2 \succeq 0$.
\end{definition}

\begin{theorem}[Spectral theorem]
Every real symmetric matrix factors as $P = U \Lambda U\T$ with $U$ orthogonal and
$\Lambda = \diag(\lambda_1,\dots,\lambda_n)$ real. $P \succeq 0$ if and only if all
$\lambda_i \ge 0$, and $P \succ 0$ if and only if all $\lambda_i > 0$.
\end{theorem}

The geometric picture is an ellipsoid. For $P \succ 0$, the set $\{x : x\T P\inv x \le 1\}$ is an
ellipsoid with principal axes along the columns of $U$ and semi-axis lengths $\sqrt{\lambda_i}$.
This is the ``one-sigma'' uncertainty ellipsoid of a Gaussian with covariance $P$
(\cref{ch:prob}).

\begin{proposition}[Properties of PSD matrices]
\label{prop:psd}
Let $P, Q \succeq 0$ and let $A$ be any matrix of compatible size. Then
\begin{enumerate}
  \item $A P A\T \succeq 0$ (congruence preserves semi-definiteness);
  \item $P + Q \succeq 0$, and $P + Q \succ 0$ if either is PD;
  \item $P_{ii} \ge 0$ and $|P_{ij}| \le \sqrt{P_{ii} P_{jj}}$;
  \item every principal submatrix of $P$ is PSD.
\end{enumerate}
\end{proposition}
\begin{proof}
(1) $x\T A P A\T x = (A\T x)\T P (A\T x) \ge 0$. (2) is immediate. (3) and (4) follow from (1) with
$A$ a selection matrix, using the $2\times2$ determinant for the off-diagonal bound.
\end{proof}

Property (1) explains why the covariance updates $\Phi P \Phi\T + Q$ and the Joseph form
(\cref{ch:kf}) preserve semi-definiteness exactly in exact arithmetic: they are sums of congruences
of PSD matrices. The simpler update $(I - KH)P$ is \emph{not} a congruence, which is why it can lose
positive semi-definiteness under rounding.

% ------------------------------------------------------------------------------
\section{Matrix factorizations}

\subsection{Cholesky and LDL$^\mathsf{T}$}

\begin{theorem}[Cholesky factorization]
A symmetric matrix $S$ is positive definite if and only if $S = L L\T$ for a unique lower
triangular $L$ with positive diagonal.
\end{theorem}

The factor is computed column by column:
\begin{equation}
  L_{jj} = \Big( S_{jj} - \sum_{k<j} L_{jk}^2 \Big)^{1/2}, \qquad
  L_{ij} = \frac{1}{L_{jj}}\Big( S_{ij} - \sum_{k<j} L_{ik}L_{jk} \Big), \quad i > j,
\end{equation}
at a cost of about $n^3/3$ flops, half that of an LU factorization. The algorithm is backward
stable without pivoting, and it fails (a negative argument of the square root) exactly when $S$ is
not numerically positive definite. That makes it a useful diagnostic too. The variant
$S = L D L\T$ with unit-diagonal $L$ and diagonal $D$ avoids square roots, and with symmetric
pivoting (Eigen's \code{LDLT}) it also handles semi-definite matrices gracefully.

Cholesky factors \emph{whiten} correlated quantities. If $n \sim \N(0, S)$ and $S = LL\T$, then
$L\inv n \sim \N(0, I)$. This is how Mahalanobis distances are computed stably:
$r\T S\inv r = \|L\inv r\|^2$, with one triangular solve.

\subsection{QR factorization and Householder reflections}

\begin{definition}[Householder reflection]
For a nonzero $v \in \R^m$, the matrix $H_v = I - 2\, \dfrac{v v\T}{v\T v}$ reflects vectors across
the hyperplane orthogonal to $v$. It is symmetric and orthogonal, $H_v = H_v\T = H_v\inv$.
\end{definition}

Given any $x \in \R^m$, choosing $v = x + \operatorname{sign}(x_1)\|x\|\,e_1$ gives
$H_v x = -\operatorname{sign}(x_1)\|x\|\,e_1$: a single reflection zeroes every entry of $x$ below
the first. (The sign choice avoids cancellation when $x$ is already nearly parallel to $e_1$.)
Applying such reflections to the columns of $A \in \R^{m\times n}$, $m \ge n$, one after another
produces
\begin{equation}
  H_n \cdots H_2 H_1\, A = \begin{bmatrix} R_1 \\ 0 \end{bmatrix}
  \quad\Longleftrightarrow\quad
  A = Q \begin{bmatrix} R_1 \\ 0 \end{bmatrix} = \begin{bmatrix} Q_1 & Q_2 \end{bmatrix}
      \begin{bmatrix} R_1 \\ 0 \end{bmatrix} = Q_1 R_1,
  \label{eq:qr}
\end{equation}
with $Q = H_1 H_2\cdots H_n$ orthogonal, $R_1 \in \R^{n\times n}$ upper triangular, $Q_1 \in
\R^{m\times n}$ and $Q_2 \in \R^{m\times(m-n)}$. The cost is about $2mn^2 - \tfrac23 n^3$ flops.

\begin{proposition}[Nullspaces from QR]
\label{prop:qr-nullspace}
If $A$ has full column rank $n$ and $A = Q_1 R_1$ as in \eqref{eq:qr}, then the columns of $Q_1$
are an orthonormal basis of $\range(A)$ and the columns of $Q_2$ are an orthonormal basis of the
left nullspace $\nullsp(A\T)$. In particular $Q_2\T A = 0$.
\end{proposition}
\begin{proof}
$Q\T A = [R_1;\ 0]$, so $Q_2\T A = 0$. The columns of $Q_2$ are orthonormal and there are $m - n$ of
them, which is the dimension of $\nullsp(A\T)$.
\end{proof}

\begin{implnote}[never form $Q$]
$Q$ is $m\times m$ and dense, but it is defined by only $n$ reflection vectors. Libraries store
those vectors compactly in the lower triangle of the factored matrix. To compute $Q\T B$ for some
other matrix $B$, apply the $n$ reflections to $B$ one at a time, at a cost of $\mathcal{O}(mnk)$
for $B \in \R^{m\times k}$. Forming $Q$ explicitly would cost $\mathcal{O}(m^2 n)$ and $m^2$ storage.
In Eigen, \code{HouseholderQR::householderQ()} returns a lazy \code{HouseholderSequence} that can
be applied to $B$ without materializing $Q$. Whether a particular call allocates must be checked
with the allocation guard, not assumed.
\end{implnote}

A \emph{Givens rotation} zeroes a single entry by rotating two rows. It is the tool of choice when
the matrix is already nearly triangular or very sparse. The original MSCKF paper suggests Givens
rotations for its nullspace projection, because the feature Jacobian has a sparse block structure.

\subsection{The singular value decomposition}

\begin{theorem}[SVD]
Every $A \in \R^{m\times n}$ factors as $A = U \Sigma V\T$ with $U, V$ orthogonal and $\Sigma$
diagonal with entries $\sigma_1 \ge \sigma_2 \ge \dots \ge 0$, the \emph{singular values}.
\end{theorem}

The SVD reveals everything about $A$ at once. The rank is the number of nonzero singular values; the
first $r$ columns of $U$ and $V$ span the column and row spaces, and the remaining columns span the
left nullspace and the nullspace. The \emph{condition number} $\kappa(A) = \sigma_1/\sigma_n$
measures how much relative errors in the data can be amplified in the solution of a linear system.
The minimum-norm least-squares solution of $Ax \approx b$ is $x = A^+ b$ with the pseudo-inverse
$A^+ = V \Sigma^+ U\T$. The SVD is the most robust of the factorizations and also the most
expensive. In this project it appears in the linear triangulation of \cref{ch:multiview} and in
diagnostics.

% ------------------------------------------------------------------------------
\section{Solving linear systems and least squares}

\begin{pitfall}[explicit inverses]
Formulas such as $K = P H\T S\inv$ are notation, not algorithms. Computing $S\inv$ and then
multiplying costs about three times as much as solving, is less accurate, and destroys structure.
Factor $S$ once (Cholesky or LDL$^\mathsf{T}$) and solve: $S K\T = H P\T$ gives $K\T$ by two
triangular substitutions per column. Writing \code{S.ldlt().solve(Identity)} is still an explicit
inverse in disguise.
\end{pitfall}

\paragraph{Least squares.} Given an overdetermined system $Ax \approx b$ with $A \in \R^{m\times n}$
of full column rank, the least-squares solution minimizes $\|Ax - b\|^2$. Setting the gradient to zero
gives the \emph{normal equations} $A\T A x = A\T b$. They are cheap but square the condition number,
$\kappa(A\T A) = \kappa(A)^2$, which can lose half the significant digits. The QR approach is more
stable: with $A = Q_1R_1$, $\|Ax - b\|^2 = \|R_1 x - Q_1\T b\|^2 + \|Q_2\T b\|^2$, so
$x = R_1\inv Q_1\T b$ by back substitution, and the residual norm is $\|Q_2\T b\|$.

\paragraph{Weighted least squares.} If the errors have covariance $W$, minimize
$(Ax - b)\T W\inv (Ax - b)$. With $W = LL\T$ this is ordinary least squares on the whitened system
$L\inv A x \approx L\inv b$.

% ------------------------------------------------------------------------------
\section{Block matrices and the Schur complement}

Estimation constantly partitions vectors and matrices into blocks: IMU state versus clones, poses
versus landmarks. Block elimination is the basic tool.

\begin{proposition}[Block factorization and inverse]
\label{prop:block-inverse}
Let $M = \begin{bmatrix} A & B \\ C & D \end{bmatrix}$ with $A$ invertible, and define the
\emph{Schur complement} of $A$ in $M$, $M/A := D - C A\inv B$. Then
\begin{equation}
  M = \begin{bmatrix} I & 0 \\ C A\inv & I \end{bmatrix}
      \begin{bmatrix} A & 0 \\ 0 & M/A \end{bmatrix}
      \begin{bmatrix} I & A\inv B \\ 0 & I \end{bmatrix},
  \qquad \det M = \det A \cdot \det(M/A),
\end{equation}
and if $M/A$ is invertible,
\begin{equation}
  M\inv = \begin{bmatrix}
    A\inv + A\inv B\, (M/A)\inv\, C A\inv & -A\inv B\, (M/A)\inv \\
    -(M/A)\inv C A\inv & (M/A)\inv
  \end{bmatrix}.
\end{equation}
\end{proposition}
\begin{proof}
Multiply out the three factors. The inverse is the product of the inverses of the three factors in
reverse order, and the inverses of the unit-triangular block factors are obtained by negating their
off-diagonal blocks.
\end{proof}

For a symmetric $M$ with $A \succ 0$: $M \succ 0$ if and only if $M/A \succ 0$. The Schur
complement is the algebraic form of \emph{eliminating} a block of variables. \Cref{ch:prob} shows it
is exactly Gaussian marginalization in information form and conditioning in covariance form.

\begin{theorem}[Woodbury matrix identity]
\label{thm:woodbury}
For invertible $A$ and $C$ and conformable $U, V$,
\begin{equation}
  (A + U C V)\inv = A\inv - A\inv U \big( C\inv + V A\inv U \big)\inv V A\inv .
\end{equation}
\end{theorem}
\begin{proof}
Write the block matrix $\begin{bmatrix} A & U \\ V & -C\inv\end{bmatrix}$ and compute the
top-left block of its inverse in two ways with \cref{prop:block-inverse}: eliminating $A$ first gives
the right-hand side, and eliminating $-C\inv$ first gives $(A + UCV)\inv$.
\end{proof}

Woodbury connects the two forms of the Kalman update (\cref{ch:kf}). It says that a low-rank update
of a matrix gives a low-rank update of its inverse, which is why an update with $m$ measurements
needs only an $m\times m$ inverse no matter how large the state is.

% ------------------------------------------------------------------------------
\section{The matrix exponential}
\label{sec:matexp}

\begin{definition}
For $A \in \R^{n\times n}$, $\exp(A) = \sum_{k=0}^\infty \frac{A^k}{k!} = I + A + \tfrac12 A^2 + \tfrac16 A^3 + \cdots$.
The series converges absolutely for every $A$.
\end{definition}

\begin{proposition}[Properties of the matrix exponential]
\label{prop:matexp}
\begin{enumerate}
  \item $x(t) = \exp(tA)\, x_0$ is the unique solution of $\dot x = Ax$, $x(0) = x_0$.
  \item $\frac{\dd}{\dd t}\exp(tA) = A \exp(tA) = \exp(tA)\, A$.
  \item $\exp(A)\inv = \exp(-A)$, and $\det\exp(A) = e^{\tr A}$.
  \item $\exp(T A T\inv) = T \exp(A)\, T\inv$ for invertible $T$.
  \item If $AB = BA$ then $\exp(A + B) = \exp(A)\exp(B)$. \textbf{In general this fails.}
\end{enumerate}
\end{proposition}

The failure of property 5 for non-commuting matrices is the algebraic root of everything subtle
about rotations: two rotations generated by $\sx{a}$ and $\sx{b}$ do not compose into the rotation
generated by $\sx{a + b}$. \Cref{ch:lie} measures the discrepancy precisely.

\begin{example}[Nilpotent generator: the double integrator]
\label{ex:double-integrator}
For $F = \begin{bmatrix} 0 & 1 \\ 0 & 0 \end{bmatrix}$, $F^2 = 0$, so the series terminates:
$\exp(F\Dt) = \begin{bmatrix} 1 & \Dt \\ 0 & 1\end{bmatrix}$. This is the exact transition matrix of
position and velocity under zero acceleration. The inertial error dynamics of \cref{ch:errordyn}
have a similar nilpotent structure, so their exponentials can be written in closed form.
\end{example}

% ------------------------------------------------------------------------------
\section{Derivatives and Jacobians}

\begin{definition}[Jacobian layout]
For $f: \R^n \to \R^m$, the Jacobian $\frac{\partial f}{\partial x} \in \R^{m\times n}$ has
entries $\big[\frac{\partial f}{\partial x}\big]_{ij} = \frac{\partial f_i}{\partial x_j}$
(numerator layout), so that $f(x + \delta) = f(x) + \frac{\partial f}{\partial x}\delta +
o(\|\delta\|)$. Column $j$ is the sensitivity of the output to the $j$-th input direction. This
is exactly what a finite-difference test computes column by column.
\end{definition}

With this layout the chain rule reads
$\frac{\partial (g\circ f)}{\partial x} = \frac{\partial g}{\partial y}\big|_{y = f(x)}\frac{\partial f}{\partial x}$,
a product of matrices in the same order as the composition. A few derivatives recur:
\begin{equation}
  \frac{\partial (Ax)}{\partial x} = A, \qquad
  \frac{\partial (a \times x)}{\partial x} = \sx{a}, \qquad
  \frac{\partial (x \times a)}{\partial x} = -\sx{a}, \qquad
  \frac{\partial (x\T S x)}{\partial x} = x\T (S + S\T).
\end{equation}
\Cref{ch:lie} extends Jacobians to functions whose inputs or outputs are rotations.

% ------------------------------------------------------------------------------
\section{Floating-point arithmetic}
\label{sec:floating}

An IEEE~754 \code{double} stores a sign, an 11-bit exponent, and a 53-bit significand (52 stored
bits plus an implicit leading one). Two constants matter:
\begin{itemize}
  \item the \emph{machine epsilon} $\mach = 2^{-52} \approx 2.22\times10^{-16}$, the gap between $1$
        and the next larger double (\code{std::numeric\_limits<double>::epsilon()});
  \item the \emph{unit roundoff} $u = 2^{-53} \approx 1.11\times10^{-16}$: every basic operation
        satisfies $\mathrm{fl}(a \circ b) = (a\circ b)(1 + \delta)$ with $|\delta| \le u$.
\end{itemize}
Relative accuracy is therefore excellent for a single operation. Trouble comes from
\emph{cancellation}: subtracting two nearly equal numbers that each carry a rounding error leaves
a small result whose relative error can be enormous.

\begin{example}[Catastrophic cancellation in $1 - \cos\theta$]
\label{ex:cancellation}
For $\theta = 10^{-5}$, $\cos\theta = 1 - 5\times10^{-11} + \dots$. The stored value of
$\cos\theta$ has an absolute error up to $u \approx 10^{-16}$, so $1 - \cos\theta \approx
5\times 10^{-11}$ is computed with relative error up to $10^{-16}/(5\times10^{-11}) \approx 2\times
10^{-6}$. Dividing by $\theta^2$ to form $(1-\cos\theta)/\theta^2 \approx 1/2$ keeps that relative
error. At $\theta = 10^{-8}$, $\cos\theta$ rounds to exactly $1$, and the computed value is $0$
instead of $1/2$. The Taylor expansion $\tfrac12 - \tfrac{\theta^2}{24} + \cdots$ has no
cancellation at all. \Cref{fig:taylor} in \cref{ch:lie} shows this effect for all the coefficient
functions of the rotation group.
\end{example}

\begin{pitfall}[time stamps in floating point]
EuRoC timestamps are nanoseconds since the Unix epoch, about $1.4\times10^{18}$, which lies between
$2^{60}$ and $2^{61}$. A \code{double} near that value can only represent multiples of
$2^{60-52} = 256$~ns. Converting absolute timestamps to \code{double} therefore quantizes them to
256~ns, and accumulating time by repeated addition of $\Dt$ drifts further. Keep absolute time in
64-bit integers (\code{std::chrono::nanoseconds}) and convert only differences to \code{double}
(AGENTS.md C5).
\end{pitfall}

\paragraph{Finite-difference step size.} A finite-difference derivative is itself a cancellation.
For the central difference $\frac{f(x+h) - f(x-h)}{2h}$, the truncation error is about
$\frac{h^2}{6}|f'''|$ and the rounding error about $\frac{u |f|}{h}$. Their sum is minimized at
\begin{equation}
  h^\star = \left(\frac{3\,u\,|f|}{|f'''|}\right)^{1/3} \approx 10^{-5}\text{ to }10^{-6}
  \quad\text{(for quantities of order one)},
\end{equation}
with a total error of order $u^{2/3} \approx 10^{-11}$. For the one-sided difference the optimal step
is $\sim\sqrt{u} \approx 10^{-8}$ and the error is only $\sim 10^{-8}$. This is why Jacobian tests
use central differences with $h \approx 10^{-6}$ and tolerances around $10^{-6}$ to $10^{-8}$, scaled
to the magnitude of the quantities involved.

\paragraph{Conditioning versus stability.} A problem is \emph{ill-conditioned} if small changes in
its data produce large changes in its exact answer. An algorithm is \emph{unstable} if it produces
much larger errors than the conditioning of the problem justifies. Computing the rotation angle
from a rotation matrix with $\arccos$ near zero angle is an unstable algorithm for a
well-conditioned problem (\cref{ch:lie}). Triangulating a point seen with almost no parallax is an
ill-conditioned problem that no algorithm can rescue (\cref{ch:multiview}).

\begin{takeaways}
\begin{itemize}
  \item $\sx{a}b = a\times b = -\sx{b}a$, $\sx{a}^3 = -\|a\|^2\sx{a}$, and $\sx{Ra} = R\sx{a}R\T$ are
        the workhorses of every rotation derivation.
  \item The left nullspace of a tall matrix is the orthogonal complement of its column space.
        A full QR factorization delivers an orthonormal basis for it.
  \item Covariance updates should be congruences ($APA\T$) plus PSD terms, so that positive
        semi-definiteness is preserved structurally.
  \item Factor and solve; never form an explicit inverse, and never form $Q$ explicitly.
  \item Schur complements eliminate variables; Woodbury converts between large and small inverses.
  \item Cancellation, not single operations, destroys floating-point accuracy. Time stamps belong
        in 64-bit integers.
\end{itemize}
\end{takeaways}

\begin{exercises}
  \item Prove \eqref{eq:skew-commutator} directly from \eqref{eq:skew-prod}, and verify the Jacobi
        identity numerically for random vectors.
  \item Show that the eigenvalues of $\sx{a}$ are $0$ and $\pm i\|a\|$. Which eigenvector belongs to $0$?
  \item Let $H_f \in \R^{2M\times3}$ have full column rank. Compute the dimension of its left
        nullspace for $M = 2, 3, 10$. Prove \eqref{eq:nullspace-projector}.
  \item Show that if $P\succ0$ then $P - K S K\T \succeq 0$ for $S = HPH\T + R$, $K = PH\T S\inv$
        and $R \succ 0$. (Hint: Schur complement of a suitable $2\times2$ block matrix.)
  \item Implement Householder QR in Python for a $10\times3$ matrix and check
        \cref{prop:qr-nullspace} numerically. Then compare the accuracy of solving a least-squares
        problem with $\kappa(A) = 10^8$ by the normal equations and by QR.
  \item Use \cref{thm:woodbury} to show that
        $\big(P\inv + H\T R\inv H\big)\inv = P - PH\T(HPH\T + R)\inv HP$. Which side is cheaper to
        evaluate when $P$ is $75\times75$ and $R$ is $2\times2$? When $R$ is $500\times500$?
  \item What is the spacing of doubles near $2^{30}$ seconds expressed in nanoseconds? Would storing
        time as \code{double} seconds since the epoch be acceptable for a 200~Hz IMU?
  \item For $f(x) = \sin x$ at $x = 1$, plot the error of forward and central differences versus
        $h \in [10^{-16}, 10^{-1}]$ on log axes and locate the optimal step sizes. Compare with the
        formulas above.
\end{exercises}
